\chapter{Introduction}

This chapter introduces the background of self-organized criticality (SOC) and the Bak--Tang--Wiesenfeld (BTW) sandpile model. It also explains the theoretical importance of the Abelian property and the main aim of this project. 

\section{Complex Systems and Self-Organized Criticality}

In many physical systems, the size of an effect is often related to the size of its cause. For example, there may only be a local impact when a small meteorite hits the Earth, but a much larger object may influence the whole orbital system. However, in complex system, this simple relationship does not always hold. In some complex systems, a small local perturbation can lead to a large-scale event \citep{Ladyman2013}.

A complex system usually contains many components which could interact with each other. Since these components affect each other, the global behaviour of the system is difficult to predict only from the local rules \citep{Ladyman2013}. And some systems can produce events with very different sizes. For example, a small change of stress may cause a large earthquake; a small ignition source may cause a large forest fire; a small fluctuation in a financial market may also spread and lead to a wider market crash. 

Bak, Tang, and Wiesenfeld introduced the idea of self-organized criticality (SOC) in 1987 \citep{Bak1987}. The SOC describes systems which can automatically evolve to a critical state, without tuning external parameters. In the critical state, the system can produce both small and large events. The system also often shows scale-invariant behaviour \citep{TADIC2023}. 

Scale invariance means that, at different scales, the statistical behaviour of the system is similar \citep{TADIC2023}. And scale invariance is closely related to power-law behaviour \citep{Aschwanden2025}. This means that small events happen very frequently, while large events happen less frequently but they are still possible. For example, this type of behaviour is commonly observed in earthquakes, forest fires, and financial market fluctuations \citep{Aschwanden2025} \citep{TADIC2023}. 

Mathematically, a power-law distribution can arise from scale-type variable transformations \citep{Sornette2006}. If the probability measure is preserved under a change of variable,
$$
P_X(x)\,dx = P_Y(y)\,dy,
$$
then the transformed distribution depends on the Jacobian of the transformation.
For a scale-type transformation such as
$$
y=x^{-1/\alpha}, \qquad \alpha>0,
$$
the transformed probability density becomes
$$
P_Y(y)=P_X(x)\left|\frac{dx}{dy}\right|.
$$
If \(P_X(x)\) is approximately constant near the origin, this gives
$$
P_Y(y)\propto y^{-(1+\alpha)},
$$
which is a power-law distribution.

This shows that scale-related transformations can naturally produce power-law behaviour. Therefore, power-law distributions are a natural description for systems with scale-invariant behaviour. (\textbf{Ans. to Q1})

For studying SOC, the Bak--Tang--Wiesenfeld (BTW) sandpile model is a simple model to begin with. Although its rules are simple, the sandpile model can produce complex avalanche behaviour. By adding grains of sand slowly and applying local toppling rules, the model can reach a critical state and show statistics with power-law \citep{Bak1987}. Therefore, the BTW model is useful for us to study complex systems, self-organized criticality, and non-equilibrium statistical behaviour.



\section{The Bak–Tang–Wiesenfeld (BTW) sandpile model}

The Bak--Tang--Wiesenfeld (BTW) sandpile model simulates a real sandpile on a finite two-dimensional lattice. Although the real sandpile is not created exactly on a the grid, this simplified model keeps the main mechanisms of the system, which include slow input, local instability, chain reactions, and boundary loss \citep{Bak1987}. 

In the BTW model, the sandpile is performed on an $N \times M$ grid. Each site $(i,j)$ has a height $h(i,j)$ as an integer, which represents the number of sand grains at that site. At each time step, one grain of sand is added to one site. This process represents a slow and local external driving force \citep{Watkins2016}.

When the height of a site reaches or exceeds the given stability threshold, the site becomes unstable. In the standard two-dimensional BTW model, the threshold is usually denoted as 4. If a site satisfies: 
$$
h(i,j) \geq 4,
$$
then this site topples. 

During a toppling process, the site loses 4 grains of sand, and these grains are assigned to its four orthogonal neighbours: up, down, left, and right \citep{Bak1987}. This rule can be mathematically written as
$$
h(i,j) \leftarrow h(i,j) - 4,
$$
$$
h(i\pm1,j) \leftarrow h(i\pm1,j) + 1,
$$
$$
h(i,j\pm1) \leftarrow h(i,j\pm1) + 1.
$$

If a toppling occurs at the grid boundary, some grains may be sent outside the lattice, and these grains will be removed from the system and counted as boundary loss. The boundary loss is important, because it prevents the total amount of sand on the grid from increasing without limit. With continuous sand addition and boundary loss, the system can eventually reach a statistically stable state \citep{Bak1987} \citep{Watkins2016}. 

The toppling of one site may cause an unstable condition for its neighbours, and these kind of sites may topple and create further unstable sites. Then a chain reaction of topplings happens, and this event caused by adding one grain of sand is called an avalanche \citep{Watkins2016}. The behaviour of avalanches is the main object studied in this project, because is shows how a small local input can create events with different sizes. 

According to above introduction for the BYW model, we can observe several important properties required for a system to exhibit SOC: First, the BTW model is a slowly driven system, we only add one grain of sand to the lattice at each time step; Second, the model contains a local threshold mechanism, since the site topples when the height of a site reaches 4 or more; Third, the BTW model can naturally produces avalanche behaviour, since one toppling event may cause neighbouring sites to become unstable and then create a chain reaction of more topplings. Fourth, the model includes boundary dissipation, since some grains leave the system when toppling occurs at the edge of the lattice; Fifth, the BTW model can evolve toward a critical state without carefully tuning parameters \citep{Bak1987}; Finally, the BTW model exhibits scale invariance and power-law behaviour \citep{Bak1987}. (\textbf{Ans. to Q2})

In the BTW model, finite boundaries are important because they allow sand grains to leave the system when toppling occurs. This boundary dissipation can prevent the total amount of sand from increasing without limit. As a result, the system can reach a statistically stable critical state under continuous sand addition. However, if we consider infinite lattice and drop sand on only a finite region, grains could spread arbitrarily far away. (\textbf{Ans. to Q3})

In the simulation in this study, four avalanche properties are recorded: topples, area, loss, and length. The value of topples is the total number of toppling events in one avalanche; The value of area is the number of unique sites that topple at least once. The value of loss is the number of grains that leave the lattice through the boundary. The value of length measures the spatial extent of the avalanche, in this study, length is defined as the avalanche radius. 

By implementing above sand addition, toppling, and avalanche simulation, we can study the behaviour of the system, the distributions of avalanche properties, the correlations between variables. 



\section{Abelian Property and Theoretical Importance}

Besides studying avalanche through simulation, the BTW sandpile model also has an important theoretical property, which is called the Abelian property. This property means that, during one avalanche, the order of topplings does not influence the final relaxed configuration. If there are several unstable sites at the same time, these sites can be toppled in different orders. Dhar states that toppling site $i$ before site $i'$ gives the same configuration as toppling site $i'$ before site $i$, so the toppling operations commute \citep{DHAR1999}. Therefore, when all necessary topplings are completed, the final stable sandpile configuration does not depend on the chosen toppling order.

This property is important for the simulation. In one avalanche, there may be many topplings, and there may be more than one unstable site at the same time. If the model did not have the Abelian property, the code would need to control the exact order of topplings very carefully. This would make the simulation more difficult to implement. However, because the BTW model has the Abelian property, the avalanche can be processed using a simpler algorithm.

The Abelian property is also important from a theoretical aspect. It shows that the BTW model has a clear mathematical structure, even though its local rules are simple. In the project document, this property is introduced by using the toppling matrix, the toppling operator, and the stabilization operator. The toppling matrix describes how the height variables are updated when a site topples \citep{DHAR1999}. In our project notation, the toppling operator represents this local update from one configuration to another. The stabilization operator then represents the full relaxation process, where an unstable configuration is turned into a stable configuration through a sequence of topplings \citep{DHAR1999}.

Therefore, the abelian property directly supports the code design and the reliability of the simulation results in this project. This allows the simulation to record avalanche properties such as topples, area, loss, and length in a consistent way. 



\section{Aim of The Project}

The main aim of this project is to implement a Bak--Tang--Wiesenfeld (BTW) sandpile model, and use it to study self-organized criticality and the related statistics on a finite two-dimensional lattice. At the end, we will extend the model to investigate different system behaviours. 